\subsection{性质（GDF）}

在本目中，我们将考察具有以下性质（称为“可数闭图”性质）的局部凸空间 F ： \(^{6}\)

(GDF) 若 u 是 F 到 Banach 空间 B 中的线性映射，使得在乘积空间 \(F \times B\) 中， u 的图像 \(\Gamma\) 的点的每个收敛序列的每个极限仍属于 \(\Gamma\) ，则 u 连续。

每个 Fréchet 空间都具有性质（GDF）（TVS, I, §3, No. 3, 定理 1 的推论 5）。在附录中，我们将看到具有性质（GDF）的空间的其他例子。

命题 11. --- 每个具有性质（GDF）的 Hausdorff 局部凸空间 F 都是桶式的。

令 V 为 F 中的一个桶， q 是它的度规，它是 F 上的一个半范数；令 H 为与赋以此单个半范数所定义的拓扑的空间 F 相关联的 Hausdorff 空间。H 的完备化 \(\hat{H}\) 是一个 Banach 空间；令 \(\pi\) 为 F 到 \(\hat{H}\) 中的典范映射；我们将证明 \(\pi\) 是连续的（对 F 上的原拓扑）；这样命题就得到证明，因为 V —— \(\pi\) 下 \(\hat{H}\) 的单位球的逆像——届时将是 F 中 0 的一个邻域。为建立 \(\pi\) 的连续性，根据（GDF），只需证明 \(\pi\) 的图像在 \(F \times \hat{H}\) 中是闭的；换言之，我们必须看到：若 \(\mathfrak{F}\) 是 F 上的一个滤子，收敛于 \(x \in F\) ，且它的像 \(\pi(\mathfrak{F})\) 收敛于 \(y \in \hat{H}\) ，则 \(y = \pi(x)\) 。而今，元素 \(x'\) （跑遍极集 \(V^{\circ}\) ，即 V 在 \(F'\) 中的极集）都可以以唯一的方式延拓为 \(\hat{H}\) 上的连续线性形式（仍记作 \(x'\) ），且这些形式所成之集是 \(\hat{H}\) 的对偶的单位球；因此只需证明 \(\langle y, x' \rangle = \langle \pi(x), x' \rangle\) （对每个 \(x' \in V^{\circ}\) ）。而这由关系式

\[
\langle \mathbf {y}, \mathbf {x} ^ {\prime} \rangle = \lim _ {\mathfrak {F}} \left\langle \pi (\mathbf {z}), \mathbf {x} ^ {\prime} \right\rangle = \lim _ {\mathfrak {F}} \left\langle \mathbf {z}, \mathbf {x} ^ {\prime} \right\rangle = \left\langle \mathbf {x}, \mathbf {x} ^ {\prime} \right\rangle = \left\langle \pi (\mathbf {x}), \mathbf {x} ^ {\prime} \right\rangle .
\]

定理 1 （Gelfand--Dunford）. --- 设 F 是具有性质（GDF）的 Hausdorff 局部凸空间， \(F_{s}^{\prime}\) 是它的弱对偶。对 T 到 \(F_{s}^{\prime}\) 中的每个标量本质 \(\mu\)-可积映射 f ，积分 \(\int f d\mu\) 属于 \(F^{\prime}\) 。

回忆一下， \(F_{s}^{\prime}\) 的对偶是 F（TVS, II, §6, No. 2, 命题 3）。于是对每个 \(z \in F\) ，数值函数 \(\langle z, f \rangle\) 是本质 \(\mu\)-可积的；令 \(\theta(z)\) 为它在 \(L^{1}(\mu)\) 中的类。为证明 \(\int f d\mu \in F^{\prime}\) ，必须证明线性形式 \(z \mapsto \langle z, \int f d\mu \rangle\) 在 F 上连续；实际上，我们将证明以下更强的结果：

引理 2. --- 设 f 是 T 到 \(F_{s}^{\prime}\) 中的映射，使得对每个 \(z \in F\) ，数值函数 \(\langle z, f \rangle\) 属于 \(\overline{\mathcal{L}^{p}}(\mu) (1 \leqslant p \leqslant +\infty)\) ；令 \(\theta(z)\) 为该函数在 \(L^{p}(\mu)\) 中的类。则 \(z \mapsto \theta(z)\) 是 F 到 \(L^{p}(\mu)\) 中的连续线性映射。

根据性质（GDF），只需证明：对 F 中元素的每个收敛于 z 的序列 \((\mathbf{z}_{n})\) ，若 \((\theta(\mathbf{z}_{n}))\) 收敛于 \(u \in \mathrm{L}^{p}(\mu)\) ，则有 \(u = \theta(\mathbf{z})\) 。而今，若有必要，把序列 \((\mathbf{z}_{n})\)

换成它的一个子列，就可以假设函数序列 \(\langle \mathbf{z}_n,\mathbf{f}\rangle\) 局部几乎处处收敛到一个函数 \(h\in \overline{\mathcal{L}^p} (\mu)\) ， \(u\) 是它在 \(\mathrm{L}^p (\mu)\) 中的类（Ch. IV, §3, No. 4, 定理 3 以及 Ch. V, §1, No. 3）。由于按假设，对每个 \(t\in \mathrm{T}\) ，序列 \((\langle \mathbf{z}_n,\mathbf{f}(t)\rangle)\) 收敛于 \(\langle \mathbf{z},\mathbf{f}(t)\rangle\) ，故局部几乎处处有 \(h(t) = \langle \mathbf{z},\mathbf{f}(t)\rangle\) ，从而 \(u = \theta (\mathbf{z})\) 。

推论 1. --- 设 \(G_{i}\) ( \(1 \leqslant i \leqslant n\) ) 是 n 个具有性质（GDF）的 Hausdorff 局部凸空间， F 是 \(\prod_{i=1}^{n} G_{i}\) 上分别连续的多线性形式所成的空间，赋以逐点收敛拓扑。对 T 到 F 中的每个标量本质 \(\mu\)-可积映射 f ，有 \(\int f d\mu \in F\) 。

空间 F 与张量积 \(\bigotimes_{i=1}^{n}G_{i}\) 处于对偶，而 F 上的逐点收敛拓扑不是别的，正是拓扑 \(\sigma(F,\bigotimes_{i=1}^{n}G_{i})\) 。因此代数对偶 \(F^{\prime*}\) 是 \(\prod_{i=1}^{n}G_{i}\) 上全体多线性形式所成的空间。设 \(\mathbf{z}=(\mathbf{z}_{1},\ldots,\mathbf{z}_{n})\) 是 \(\prod_{i=1}^{n}G_{i}\) 的一个元素；对每个多线性形式 \(u\in F^{\prime*}\) ，映射 \(\mathbf{x}\mapsto\mathbf{u}(\mathbf{z}_{1},\ldots,\mathbf{z}_{i-1},\mathbf{x},\mathbf{z}_{i+1},\ldots,\mathbf{z}_{n})\) 是 \(G_{i}\) 上的一个线性形式，我们把它记作 \(\lambda_{i}(\mathbf{z})(\mathbf{u})\) ；于是我们得到线性映射 \(\lambda_{i}(\mathbf{z})\) ，它把 \(F^{\prime*}\) 映入代数对偶 \(G_{i}^{*}\) （ \(G_{i}\) 的代数对偶）中，且对拓扑 \(\sigma(F^{\prime*},\bigotimes_{i=1}^{n}G_{i})\) 与 \(\sigma(G_{i}^{*},G_{i})\) 连续。说 \(u\in F\) 就意味着：对每个指标 i 与每个 \(z\in\Pi_{i=1}^{n}G_{i}\) ，有 \(\lambda_{i}(\mathbf{z})(\mathbf{u})\in G_{i}^{\prime}\) 。于是，根据第 1 目命题 1，映射 \(\lambda_{i}(\mathbf{z})\circ\mathbf{f}\) 是标量本质 \(\mu\)-可积的、T 到 \(G_{i}^{\prime}\) （赋有拓扑 \(\sigma(G_{i}^{\prime},G_{i})\) ）中的映射，并且

\[
\int \left(\lambda_ {i} (\mathbf {z}) \circ \mathbf {f}\right) d \mu = \lambda_ {i} (\mathbf {z}) \left(\int \mathbf {f} d \mu\right).
\]

根据定理 1， \(\int (\lambda_i(\mathbf{z})\circ \mathbf{f})d\mu \in \mathrm{G}_i^\prime\) （对 \(1\leqslant i\leqslant n\) ），因此 \(\int \mathbf{f}d\mu \in \mathbf{F}\)

推论 2. --- 设 G 是具有性质（GDF）的 Hausdorff 局部凸空间， H 是半自反空间，其强对偶 \(H_{b}^{\prime}\) 具有性质（GDF）（参见附录, No. 2, 命题 3）。令 F 为空间 \(\mathcal{L}_{s}(G;H)\) ；对 T 到 F 中的每个标量本质 \(\mu\)-可积映射 U ，积分 \(\int U d\mu\) 属于 F 。

由于 G 是桶式的（命题 11）， \(\mathcal{L}(\mathrm{G};\mathrm{H}) = \mathcal{L}(\mathrm{G}_{\sigma};\mathrm{H}_{\sigma})\) （TVS, IV, §1, No. 3, 命题 7）；此外，可以把 \(\mathrm{F} = \mathcal{L}_s(\mathrm{G};\mathrm{H})\) 换成空间 \(\mathcal{L}_s(\mathrm{G}_{\sigma};\mathrm{H}_{\sigma})\) ，因为这两个空间有相同的对偶 \(\mathrm{G} \otimes \mathrm{H}'\) （TVS, II, §6, No. 2, 命题 3, 以及上文第 3 目第一段）。若对每个 \(u \in \mathcal{L}(\mathrm{G};\mathrm{H}) = \mathcal{L}(\mathrm{G}_{\sigma};\mathrm{H}_{\sigma})\) ，令 \(\widetilde{u}(\mathbf{x},\mathbf{y}') = \langle u(\mathbf{x},\mathbf{y}') \rangle\) （对 \(\mathbf{x} \in \mathbf{G}\) ， \(y' \in H')\)），则线性映射 \(u \mapsto \widetilde{u}\) 是 F 到空间 \(F_1\) （由 \(G_\sigma \times H_s'\) 上分别连续双线性型组成）上的双射，其中 \(H_s'\) 表示对偶 \(H'\) （赋以弱拓扑 \(\sigma(H', H)\) ）（附录, No. 1）；此外，该映射是 \(\mathcal{L}_s(G_\sigma; H_\sigma)\) 到赋以逐点收敛拓扑的 \(F_1\) 上的同构（出处同上）。但由于按假设 H 是 \(H_b'\) 的对偶， \(F_1\) 也是 \(G \times H_b'\) 上分别连续双线性型所成的空间。于是推论 2 由推论 1 得出。

注意，推论 2 特别适用于 G 为 Banach 空间且 H 为自反 Banach 空间的情形。

